\documentclass{beamer}

\input{../helper}

\title{Introdução à Computação \\ achando erros}
\author{Alexandre Rademaker}
\date{}

\begin{document}

\begin{frame}
 \maketitle
\end{frame}



\begin{frame}[fragile,allowframebreaks]{Depuração (debug)}

  Bugs são erros que fazem com que o programa se comporte de forma
  diferente do pretendido. E a depuração é o processo de encontrar e
  corrigir esses bugs.

  Uma forma de \textbf{debug} é imprimir valores das variáveis em
  determinados pontos do programa.  

  Uma boa técnica de depuração é explicar nosso próprio código em voz
  alta para alguém (ou alguma coisa!).

  Em cada ambiente o suporte pode ser diferente. Vamos experimentar no
  \textbf{Codespaces}!

  \framebreak

  Queremos 3 blocos na tela. Qual o bug?

  \lstinputlisting[style=normalc,caption=src/buggy0.c]{src/buggy0.c}

  \framebreak

  Qual o bug? No VS Code, usar `step into'.

  \lstinputlisting[style=tinyc,caption=src/buggy1.c]{src/buggy1.c}

  \framebreak

  Lembrando do problema em \ilcode{semana-2/src/mario2.c}. Vamos
  adicionar um pouco de robustez em \ilcode{mario3.c}.

  Queremos garantir que a entrada passada pelo usuário é compatível
  com o que esperamos.

  Usamos o \ilcode{printf} para entender os valores.

  Precisamos de algumas bibliotecas adicionais para poder usar o valor
  \ilcode{true} e a função \ilcode{isdigit}.
  
  \framebreak

  \lstinputlisting[style=tinyc,caption=src/mario3.c]{src/mario3.c}


\end{frame}


\end{document}

%%% Local Variables:
%%% mode: latex
%%% TeX-master: t
%%% End:
